\section*{摘要}
面向智能家居、安防、工业感知及\DrivingExtensionDirection{}等场景对本地响应、数据隐私和离线运行的需求，本项目提出\ProjectName{}。针对模型部署中算子接口不统一、异构计算协同困难和跨系统移植成本高等问题，采用“可重构硬件平台、统一推理软件、感知通信适配”的总体设计，围绕开源\ProcessorArchitecture{}架构与国产开源软件，贯通硬件集成、驱动开发、系统移植、模型推理及应用接入。

硬件以现有\BoardModel{}为原型载体，通过\IntegrationFramework{}集成\ProcessorCore{}通用处理器、\VectorUnit{}向量单元和\AcceleratorName{}矩阵加速器，分别承担控制与标量计算、向量化运算及数据整理、受支持的矩阵乘加，并由\MemoryController{}连接板上 DDR3，协同设计缓存、地址访问和数据搬运。系统软件分别适配国产开源\RealTimeSystem{}与可裁剪的\LinuxDistribution{}，复用推理核心源码和平台接口，通过独立镜像构建、重启切换验证两类运行环境。

软件以\RuntimeName{}为核心：主机端完成模型训练与量化准备，由受限转换器检查算子、形状及量化约束，生成后端分配、执行顺序和静态内存计划；板端加载模型与输入，按计划调用标量、RVV 或 NPU 后端，回传带样本标识的结果与诊断记录。项目的研究重点在于统一算子语义与量化规则、按计算特点分配后端，以及按张量生命周期复用缓冲区。通过独立数值参考、单算子和模型逐层对照，将开发接口、执行正确性与跨系统复现联系起来，为新增算子和后续优化提供可核验的依据。

实施上，拟以\ModelDataType{}量化的\AIModel{}为首个算例，从静态输入出发，依次开展数值对照、联合硬件仿真、资源与时序检查、板级启动及两种系统下的完整推理，再扩展摄像头与传感器采集、结果上送和移动数据中心通信适配。预期形成可配置的硬件原型、可复用的算子库与运行时，以及配套部署和验证工具，服务嵌入式模型原型开发、算子优化与教学实训，减少重复适配和问题定位的工作量，为后续国产处理器平台迁移及智能感知应用提供基础。

\par\noindent\textbf{关键词：}\ProjectKeywords{}
